\chapter{Biomarker Interpretation}
\section{Acute Injury Markers}
\textit{HAVCR1} encodes KIM-1, a well-known proximal tubular injury marker. Elevated \textit{HAVCR1} expression in AKI-related cells is biologically consistent with epithelial injury and repair activation. \textit{LCN2} encodes NGAL and is associated with tubular stress, innate immune signaling, and early kidney injury. These genes provide positive biological controls for the AKI biomarker analysis.

\section{Fibrosis and CKD Markers}
CKD progression is strongly associated with fibrosis. \textit{COL1A1} encodes type I collagen and reflects extracellular matrix accumulation. \textit{FN1} encodes fibronectin, a matrix glycoprotein involved in tissue remodeling and wound repair. \textit{SERPINE1} encodes PAI-1 and is associated with profibrotic signaling and impaired matrix degradation. Increased importance of these genes in CKD-related comparisons would support a chronic remodeling interpretation.

\section{CKD Progression Biomarkers}
The CKD versus AKI comparison is designed to identify genes that separate acute injury from chronic disease. A gene highly ranked in this comparison may represent maladaptive repair, persistent inflammation, or fibrotic transition. Such candidates are important because they may help identify patients at risk of incomplete recovery after AKI.

\section{Clinical Relevance}
Clinically useful biomarkers should be measurable, reproducible, disease-specific, and linked to meaningful outcomes. Transcriptomic candidates from this study require validation at the protein level through urine, blood, or tissue assays. Genes encoding secreted or membrane-associated proteins are especially promising for translational biomarker development.
